\TieudeT{PHONG TỎA VẬT LÍ}
\subsubsection*{PHẦN I. Thí sinh trả lời câu hỏi từ câu 1 đến câu 18. Mỗi câu hỏi thí sinh chỉ chọn một phương án}
\setcounter{ex}{0}
\begin{ex}%book
Phép đo của một đại lượng vật lí
\choice
{\True là phép so sánh nó với một đại lượng cùng loại được quy ước làm đơn vị}
{là những sai sót gặp phải khi đo một đại lượng vật lí}
{là sai số gặp phải khi dụng cụ đo một đại lượng vật lí}
{là những công cụ đo các đại lượng vật lí như thước, cân...vv}
\loigiai{
}
\end{ex}
\begin{ex}%book
Trong đơn vị SI, đơn vị nào là đơn vị dẫn xuất?
\choice
{mét (m)}
{giây (s)}
{mol (mol)}
{\True mét trên giây $(m/s)$}
\loigiai{

}

\end{ex}
\begin{ex}%book
Có mấy cách để đo các đại lượng vật lí?
\choice
{1}
{\True 2}
{3}
{4}
\loigiai{

}
\end{ex}
\begin{ex}%book
Trong các phép đo dưới đây, đâu là phép đo trực tiếp?
\begin{itemize}	
\item[(1)] Dùng thước đo chiều cao.
\item[(2)] Dùng cân đo cân nặng.
\item[(3)] Dùng cân và ca đong đo khối lượng riêng của nước.
\item[(4)] Dùng đồng hồ và cột cây số đo tốc độ của người lái xe.
\end{itemize}
\choice
{\True $(1)$, $(2)$}
{$(2)$, $(4)$}
{$(1)$, $(2)$, $(4)$}
{$(2)$, $(3)$, $(4)$}
\loigiai{
}
\end{ex}
\begin{ex}%book
Một học sinh dùng thước đo chiều dài có độ chia nhỏ nhất là 1 milimet, đồng hồ bấm giây để đo quãng đường, thời gian và tốc độ chuyển động của chiếc xe đồ chơi. Phát biểu nào sau đây là \textbf{sai}?
\choice
{\True Đo tốc độ là phép đo trực tiếp}
{Đo thời gian là phép đo trực tiếp}
{Sai số dụng cụ của thước đo chiều dài là $0,5\ mm$}
{Đo quãng đường là phép đo trực tiếp}
\loigiai{}
\end{ex}
\begin{ex}%book %[1P1YB-2][Loai1]
Kết quả sai số tuyệt đối của một phép đo là $0,0809$. Số chữ số có nghĩa là
\choice
{$4$}
{$2$}
{$1$}
{\True $3$}
\loigiai{
Số 0,0809 có 3 chữ số có nghĩa.
}
\end{ex}
\begin{ex}%book %[1P1YB-2][Loai1]
Kết quả sai số tuyệt đối của một phép đo là $0,4009$. Số chữ số có nghĩa là
\choice
{\True $4$}
{$1$}
{$2$}
{$3$}
\loigiai{
Số 0,4009 có 4 chữ số có nghĩa.
}
\end{ex}
\begin{ex}%book
Sai số hệ thống là
\choice
{\True sai số có tính quy luật và được lặp lại ở tất cả các lần đo}
{sai số từ những yếu tố ngẫu nhiên bên ngoài}
{sai số xuất phát từ sai sót, phản xạ của người làm thí nghiệm}
{sai số tuyệt đối}	
\loigiai{
}
\end{ex}
\begin{ex}%book
$\Delta A_{i}=\left|\overline{A}-\Delta A_{i}\right|$ là công thức tính
\choice
{\True sai số tuyệt đối ứng với mỗi lần đo}
{sai số tuyệt đối trung bình của n lần đo}
{giá trị trung bình của đại lượng cần đo khi tiến hành phép đo nhiều lần}
{sai số tuyệt đối của phép đo}	
\loigiai{
}
\end{ex}
\begin{ex}%book
Điều nào sau đây \textbf{không} đúng khi nói về nguyên nhân gây ra sai số ngẫu nhiên trong quá trình đo một đại lượng vật lí?
\choice
{Thao tác đo không chuẩn}
{\True Dụng cụ đo không chuẩn}
{Điều kiện làm thí nghiệm không ổn định}
{Mắt người đọc không chuẩn}
\loigiai{
}
\end{ex}
\begin{ex}%book
Sai số tỉ đối của phép đo là
\choice
{$\delta x=\dfrac{\Delta \overline{x}}{\overline{x}} .  100\%$}
{\True $\delta x=\dfrac{\Delta x}{\overline{x}} .  100\%$}
{$\delta x=\dfrac{\Delta x_{d c}}{\overline{x}} .  100\%$}
{$\delta x=\dfrac{\Delta \overline{x}}{x} .  100\%$}

\loigiai{
}
\end{ex}
\begin{ex}%book
Một phép đo đại lượng vật lí $A$ thu được giá trị trung bình là $\overline{A}$, sai số của phép đo là $\triangle A$. Cách ghi đúng kết quả đo $A$ là
\choice
{$A=\overline{A}+\Delta A$}
{$A=\overline{A}-\Delta A$}
{$A=A \pm \Delta A$}
{\True $A=\overline{A} \pm \Delta A$}

\loigiai{}

\end{ex}
\begin{ex}%book
Sai số của một tích $A= B. C$ trong phép đo gián tiếp được xác định bởi công thức?
\choice
{\True $\delta A=\delta B+\delta C$}
{$\Delta A=\Delta B+\Delta C$}
{$\delta A=\Delta B+\Delta C$}
{$\Delta A=\delta B+\delta C$}

\loigiai{
}
\end{ex}
\begin{ex}%book
Khi đo chiều dài $L$ của chiếc bàn học, một học sinh viết được kết quả là $L=118 \pm 2\ cm$. Giá trị trung bình của $L$ là
\choice
{$2\ cm$}
{$120\ cm$}
{\True $118\ cm$}
{$116\ cm$}
\haidong{2}
\loigiai{
}
\haidong{7}\end{ex}
\begin{ex}%book
Bạn Bình tiến hành cân một viên bi nhỏ nhiều lần và tính được kết quả khối lượng của viên bi như sau: $m=6,03 \pm 0,04\ g$. Sai số tuyệt đối của phép đo khối lượng viên bi là
\choice
{\True $0,04\ g$}
{$6,03\ g$}
{$0,7\%$}
{$0,6\%$}
\haidong{7}
\loigiai{
}
\end{ex}
\begin{ex}%book
Một người dùng thước đo có độ chia nhỏ nhất $1\ cm$ để đo chiều cao một người khác trong ba lần được là $160\ cm$; $161\ cm$; $162\ cm$. Lấy sai số dụng cụ bằng một độ chia nhỏ nhất. Sai số tỉ đối của phép đo là
\choice
{$0,23\%$}
{$7,2\%$}
{$0,41\%$}
{\True $1,04\%$}
\haidong{4}
\loigiai{
\begin{itemize}
\item $\overline{L}=161\ cm$.
\item $\overline{\Delta A}=\dfrac{1+0+1}{3}=\dfrac{2}{3}$
\item $\Delta A=\overline{\Delta A}+ \Delta A_{DC}=\dfrac{2}{3}+1=\dfrac{5}{3}\ cm$.
\item $\delta A=\dfrac{\Delta A}{\overline{A}}=\dfrac{\dfrac{5}{3}}{161}\approx 1,04\%$
\end{itemize}		
}
\haidong{10}\end{ex}
\begin{ex}%book
Dùng một cái cân có độ chia nhỏ nhất là $0,1\ kg$ để cân $5$ lần khối lượng của một bao gạo đều bằng $20\ kg$. Lấy sai số dụng cụ bằng một độ chia nhỏ nhất. Kết quả của phép cân này là
\choice
{$20\pm 0,1\ kg$}
{$20,0\pm 0,2\ kg$}
{$20,00\pm 0,1\ kg$}
{\True $20,0\pm 0,1\ kg$}
\loigiai{
\begin{itemize}
\item Do 5 lần cân đều cho kết quả 20 kg nên sai số ngẫu nhiên $\overline{\Delta m}=0$.
\item Sai số của phép cân bằng sai số dụng cụ: $\Delta m=\Delta m_{dc}=0,1\ kg$.
\item Kết quả của phép cân: $20,0\pm 0,1\ (kg)$.
\end{itemize}
}
\haidong{10}\end{ex}
\begin{ex}%book
Để đo dòng điện một chiều có cường độ $25\ mA$ bằng đồng hồ vạn năng, người ta tiến hành các bước:
\immini{\begin{itemize}
\item[(1)] Kết nối que đo màu đỏ của đồng hồ về phía cực (+) và que đo màu đen về phía cực (-) theo chiều dòng điện trong mạch thí nghiệm. Mắc đồng hồ nối tiếp với mạch thí nghiệm.
\item[(2)] Đọc kết quả đo hiện trên đồng hồ.
\item[(3)] Tắt nguồn điện của các mạch thí nghiệm.
\item[(4)] Cắm que đo màu đen vào đầu $C O M$ và que đỏ vào dấu $(+)$.
\item[(5)] Đặt chuyển mạch của đồng hồ ở thang $D C m A$.
\item[(6)] Bật điện cho mạch thí nghiệm.
\end{itemize}
}{\includegraphics[width=3.5cm]{images/3_7}}
Thứ tự các bước tiến hành đúng là
\choice
{\True $(4)$; $(5)$; $(3)$; $(1)$; $(6)$; $(2)$}
{$(6)$; $(5)$; $(4)$; $(3)$; $(1)$; $(2)$}
{$(1)$; $(2)$; $(3)$; $(4)$; $(5)$; $(6)$}
{$(3)$; $(1)$; $(6)$; $(5)$; $(4)$; $(2)$}

\loigiai{
}
\end{ex}
































\subsubsection*{PHẦN 2. Thí sinh trả lời từ câu 1 đến câu 4. Trong mỗi ý a), b), c), d) ở mỗi câu, thí sinh chọn đúng hoặc sai.}
\setcounter{ex}{0}
\begin{ex}
Nếu một đại lượng được đo trực tiếp.
\choiceTF[t]
{\True Giá trị trung bình của đại lượng cần đo khi tiến hành phép đo nhiều lần được xác định bằng công thức $\overline{A}=\dfrac{A_1+A_2+\ldots+A_n}{n}$}
{\True Sai số dụng cụ thường được lấy bằng một nửa độ chia nhỏ nhất trên dụng cụ}
{Sai số tỉ đối của phép đo càng lớn thì phép đo càng chính xác}
{Sai số của phép đo chiều dài một đại lượng có kết quả là $0,03 \ cm$. Số chữ số có nghĩa trong kết quả của phép đo này là 3}
\loigiai{ 
\begin{itemchoice}
\itemch
\itemch
\itemch
\itemch
\end{itemchoice}
}
\end{ex}



\begin{ex}%Tailieuchuan
Hình vẽ thể hiện nhiệt kế đo nhiệt độ $t_1\left(^\circ C\right)$ và $t_2\left(^\circ C\right)$ của một dung dịch trước và sau khi đun.
\begin{center} 
\begin{tikzpicture}[draw=Mapcolor, very thick,line cap=round, line join=round, scale=1, font=\footnotesize, >=stealth,draw=Mapcolor]

\draw[fill=cyan!30, cyan!30]  (-2.75,0.25) rectangle (-1.375,0);
\draw[Mapcolor] (0.25,0) -- (-2.75,0) -- (-2.75,0.25) -- (0.25,0.25);
\draw[Mapcolor] (-2.5,0.75) -- (-2.5,0.25);
\draw[Mapcolor] (-2.25,0.5) -- (-2.25,0.25);
\draw[Mapcolor] (-2,0.5) -- (-2,0.25);
\draw[Mapcolor] (-1.75,0.5) -- (-1.75,0.25);
\draw[Mapcolor] (-1.5,0.5) -- (-1.5,0.25);
\draw[Mapcolor] (-1.25,0.75) -- (-1.25,0.25);
\draw[Mapcolor] (-1,0.5) -- (-1,0.25);
\draw[Mapcolor] (-0.75,0.5) -- (-0.75,0.25);
\draw[Mapcolor] (-0.5,0.5) -- (-0.5,0.25);
\draw[Mapcolor] (-0.25,0.5) -- (-0.25,0.25);
\draw[Mapcolor] (0,0.75) -- (0,0.25);
\node at (-2.5,1) {$20$};
\node at (-1.25,1) {$25$};
\node at (0,1) {$30$};

\begin{scope}[shift={(5,0)},rotate=0]
\draw[fill=cyan!30, cyan!30]  (-2.75,0.25) rectangle (-0.375,0);
\draw[Mapcolor] (0.25,0) -- (-2.75,0) -- (-2.75,0.25) -- (0.25,0.25);
\draw[Mapcolor] (-2.5,0.75) -- (-2.5,0.25);
\draw[Mapcolor] (-2.25,0.5) -- (-2.25,0.25);
\draw[Mapcolor] (-2,0.5) -- (-2,0.25);
\draw[Mapcolor] (-1.75,0.5) -- (-1.75,0.25);
\draw[Mapcolor] (-1.5,0.5) -- (-1.5,0.25);
\draw[Mapcolor] (-1.25,0.75) -- (-1.25,0.25);
\draw[Mapcolor] (-1,0.5) -- (-1,0.25);
\draw[Mapcolor] (-0.75,0.5) -- (-0.75,0.25);
\draw[Mapcolor] (-0.5,0.5) -- (-0.5,0.25);
\draw[Mapcolor] (-0.25,0.5) -- (-0.25,0.25);
\draw[Mapcolor] (0,0.75) -- (0,0.25);
\node at (-2.5,1) {$60$};
\node at (-1.25,1) {$65$};
\node at (0,1) {$70$};
\end{scope}

\node at (-1.25,-0.5) {a)};
\node at (3.75,-0.5) {b)};
\end{tikzpicture}
\end{center} 
\choiceTF[t]
{\True Giá trị đo nhiệt độ $t_1\left({}^\circ C\right)$ của một dung dịch trước khi đun là\linebreak $t_1=24,0 \pm 0,5^\circ C$} 
{Giá trị đo nhiệt độ $t_1\left({}^\circ C\right)$ của một dung dịch sau khi đun là \linebreak $t_2=68,0 \pm 0,05^\circ C$} 
{Sai số tuyệt đối của phép đo nhiệt độ này là $1,5^\circ C$} 
{\True Kết quả độ tăng nhiệt độ của dung dịch là: $t=44,0 \pm 1,0^\circ C$}
\loigiai{ 
\begin{itemchoice}
\itemch
\itemch
\itemch
\itemch
\end{itemchoice}
}
\end{ex}
\begin{ex}%Tailieuchuan
Hai người cùng đo chiều dài của cánh cửa sổ, kết quả thu được như sau:
\begin{itemize}
\item Người thứ nhất: $d=120 \pm 1 \ cm$
\item Người thứ hai: $d=120 \pm 2 \ cm$
\end{itemize}
\choiceTF[t]
{\True Sai số tỉ đối được xác định bằng tỉ số giữa hai số tuyệt đối và giá trị trung bình của chiều dài cánh cửa số: $\delta d=\dfrac{\Delta d}{\overline{d}} .  100 \%$} 
{Sai số tỉ đối của phép đo của người thứ nhất là $8,3\%$} 
{\True Sai số tỉ đối của phép đo của người thứ hai là $1,67\%$} 
{Người thứ hai đo chính xác hơn người thứ nhất vì sai số tỉ đối của người thứ nhất lớn hơn} 
\loigiai{ 
\begin{itemchoice}
\itemch
\itemch
\itemch
\itemch
\end{itemchoice}
}
\end{ex}
\begin{ex}%Tailieuchuan
Một học sinh đo tốc độ trung bình của viên bi được giá trị $v=(2,50 \pm 0,04)\ m/s$.
\choiceTF[t]
{\True Giá trị trung bình của tốc độ của viên bi là $2,54 \ m/s$} 
{\True Sai số tuyệt đối của phép đo tốc độ viên bi là $0,04$} 
{\True Sai số tỉ đối của phép đo là $1,6 \%$} 
{Nếu sai số tỉ đối của thời gian là $1 \%$ thì sai số tỉ đối của quãng đường là $1,6 \%$} 
\loigiai{ 
\begin{itemchoice}
\itemch
\itemch
\itemch
\itemch
\end{itemchoice}
}
\end{ex}


\subsubsection*{PHẦN 3. Thí sinh trả lời từ câu 1 đến câu 6.}
\setcounter{ex}{0}
\begin{ex}%Tailieuchuan
Lấy sai số dụng cụ bằng nửa độ chia nhỏ nhất. Sai số dụng cụ của cây bút chì trong trường hợp dưới đây là bao nhiêu $cm$? 
\begin{center} 
\includegraphics[width=10cm]{images/2025_04_24_0674a062b382e4636f6cg-09}
\end{center} 
\shortans[oly]{0,05}
\loigiai{
}
\end{ex}
\begin{ex}
Đường kính của một quả bóng bằng $(5,0 \pm 0,2)\ cm$. Sai số tỉ đối của phép đo thể tích quả bóng có giá trị là $x\%$. Tìm x.
\shortans[oly]{12}
\loigiai{
\begin{itemize} 
\item Thể tích quả bóng là $V=\dfrac{4}{3} \pi \mathrm{R}^{3}$
\item Sai số tỉ đối của phép đo thể tích:
\begin{align*}
\delta(V)=\dfrac{\Delta V}{V} 100\%=3\dfrac{\Delta R}{R}.100 \%=3\dfrac{\Delta d}{d}.100 \%=3\dfrac{0,2}{5,0} 100 \%=12 \%
\end{align*}
\end{itemize}
}
\end{ex}
\begin{ex}%Tailieuchuan
Để tính gia tốc rơi tự do $g$, người ta có thể dùng công thức tính chu kì của một con lắc đơn gồm một vật nặng có kích thước nhỏ treo vào một dây nhẹ, không co giãn: $T=2 \pi \sqrt{\dfrac{\ell}{g}}$, trong đó, $T$ là thời gian để vật nặng thực hiện một dao động và $\ell$ là chiều dài sợi dây. Trong thí nghiệm với con lắc đơn, người ta đo được: $\ell=0,350 \pm 0,005\ (m)$ và $T=1,18 \pm 0,02\ (s)$. Sai số tuyệt đối của phép đo này là bao nhiêu $m/s^2$ (làm tròn kết quả đến chữ số hàng phần mười)?
\shortans[oly]{0,5}
\loigiai{
}
\end{ex}
\begin{ex}%Tailieuchuan
Dùng thước thẳng có giới hạn đo là $20\ cm$ và độ chia nhỏ nhất là $0,5\ cm$ để đo chiều dài chiếc bút máy. Nếu chiếc bút có độ dài cỡ $15\ cm$ thì phép đo này sai số tỉ đối là bao nhiêu $\%$ (làm tròn kết quả đến chữ số hàng phần trăm)?
\shortans[oly]{1,67}
\loigiai{
}
\end{ex}
\begin{ex}%Tailieuchuan
Kết quả đo đường kính của một đường tròn là $d=(40,1 \pm 0,2)\ cm$. Sai số tuyệt đối của phép đo diện tích này bằng  $x.10^3\ cm^2$. Tìm $x$ (làm tròn kết quả đến chữ số hàng phần trăm).
\shortans[oly]{0,01}
\loigiai{
}
\end{ex}
\begin{ex}%Tailieuchuan
Kết quả của phép đo chiều dài và chiều rộng của một hình chữ nhật lần lượt là $a=85,0 \pm 0,2\ (cm)$ và $b=29,5 \pm 0,2\ (cm)$. Sai số tỉ đối của phép đo diện tích hình chữ nhật là bao nhiêu $\%$ (làm tròn kết quả đến chữ số hàng phần mười)?
\shortans[oly]{0,9}
\loigiai{
}
\end{ex}





\TieudeT{PHONG TỎA VẬT LÍ}
\subsubsection*{PHẦN I. Thí sinh trả lời câu hỏi từ câu 1 đến câu 18. Mỗi câu hỏi thí sinh chỉ chọn một phương án}
\setcounter{ex}{0}
\begin{ex}%book
Đại lượng nào là đại lượng cơ bản của hệ SI?
\choice
{\True Cường độ dòng điện}
{Hiệu điện thế}
{Công suất}
{Điện trở}
\loigiai{
}
\end{ex}
\begin{ex}%book
Đơn vị đo thời gian trong hệ thống đo lường chính thức ở nước ta hiện nay là
\choice
{tuần}
{ngày}
{giờ}
{\True giây}
\loigiai{

}
\end{ex}
\begin{ex}%book
Phép đo trực tiếp của một đại lượng vật lí là phép
\choice
{\True so sánh trực tiếp qua dụng cụ đo}
{so sánh gián tiếp qua dụng cụ đo}
{so sánh không thông qua dụng cụ đo}
{đo đại lượng thứ nhất rồi suy ra đại lượng cần đo thông qua công thức}
\loigiai{

}
\end{ex}
\begin{ex}%book
Phép đo nào sau đây \textbf{không} thể thực hiện đo trực tiếp?
\choice
{Đo thời gian}
{Đo khối lượng}
{\True Đo khối lượng riêng}
{Đo chiều dài}
\loigiai{
}
\end{ex}
\begin{ex}%book
Thiết bị nào sau đây dùng để đo điện năng tiêu thụ trong các hộ gia đình?
\choice
{Vôn kế}
{Ampe kế}
{\True Công tơ điện}
{Nhiệt kế}
\loigiai{}
\end{ex}
\begin{ex}%book %[1P1YB-2][Loai1]
Kết quả sai số tuyệt đối của một phép đo là $6,07$. Số chữ số có nghĩa là
\choice
{$4$}
{$2$}
{\True $3$}
{$1$}
\loigiai{
Số 6,07 có 3 chữ số có nghĩa.
}
\end{ex}
\begin{ex}%book
Sai số phép đo bao gồm
\choice
{sai số ngẫu nhiên và sai số đơn vị}
{\True sai số ngẫu nhiên và sai số hệ thống}
{sai số hệ thống và sai số đơn vị}
{sai số đơn vị và sai số dụng cụ}
\loigiai{
}
\end{ex}
\begin{ex}%book
Sai số hệ thống có thể được hạn chế bằng cách thường xuyên
\choice
{hiệu chỉnh dụng cụ đo, vệ sinh dụng cụ đo}
{đeo kính lúp khi đo, vệ sinh dụng cụ đo}
{\True hiệu chỉnh dụng cụ đo, sử dụng thiết bị đo có độ chính xác cao}
{sử dụng thiết bị đo có độ chính xác cao, đeo kính lúp khi đo}
\loigiai{
}
\end{ex}
\begin{ex}%book
Khi đo n lần cùng một đại lượng $A$, ta nhận được các giá trị khác nhau: $A_1$, $A_2$, $\ldots$, $A_{n}$. Giá trị trung bình của $A$ là $\overline{A}$. Sai số tuyệt đối ứng với lần đo thứ $4$ được tính bằng công thức nào sau đây?
\choice
{\True $\Delta A_4=\left|\overline{A}-A_4\right|$}
{$\Delta A_4=\dfrac{\left|\overline{A}-A_4\right|}{2}$}
{$\Delta A_4=\dfrac{\left|\overline{A}+A_4\right|}{2}$}
{$\Delta A_4=\left|\overline{A}+A_4\right|$}
\loigiai{
}
\end{ex}
\begin{ex}%book
Sai số ngẫu nhiên có thể được hạn chế bằng cách
\choice
{\True thực hiện phép đo nhiều lần và lấy giá trị trung bình để hạn chế sự phân tán của số liệu đo}
{hiệu chỉnh dụng cụ đo, sử dụng thiết bị đo có độ chính xác cao}
{hiệu chỉnh dụng cụ đo, thực hiện phép đo nhiều lần}
{thực hiện phép đo nhiều lần, sử dụng thiết bị đo có độ chính xác cao}
\loigiai{
}
\end{ex}
\begin{ex}%book
Sai số tỉ đối của phép đo là
\choice
{tỉ số giữa sai số tuyệt đối và sai số ngẫu nhiên}
{tỉ số giữa sai ngẫu nhiên và sai số hệ thống}
{\True tỉ số giữa sai số tuyệt đối và giá trị trung bình của đại lượng cần đo}
{tỉ số giữa sai số ngẫu nhiên và sai số tuyệt đối}
\loigiai{
}
\end{ex}
\begin{ex}%book
Khi đo nhiều lần thời gian chuyển động của một viên bi trên mặt phẳng nghiêng mà thu được nhiều giá trị khác nhau, thì giá trị nào sau đây được lấy làm kết quả của phép đo?
\choice
{Giá trị trung bình của giá trị lớn nhất và giá trị nhỏ nhất}
{\True Giá trị trung bình của tất cả các giá trị đo được}
{Giá trị của lần đo cuối cùng}
{Giá trị được lặp lại nhiều lần nhất}

\loigiai{
}
\end{ex}
\begin{ex}%book
Trong bài thực hành đo tốc độ của vật chuyển động. Để đo tốc độ trung bình của vật ta sử dụng biểu thức $v=\dfrac{s}{t}$. Khi đó
\choice
{$\Delta v=\Delta s-\Delta t$}
{$\Delta v=\Delta s+\Delta t$}
{\True $\delta v=\delta s+\delta t$}
{$\delta v=\delta s-\delta t$}

\loigiai{
}
\end{ex}
\begin{ex}%book
Đo chiều dày của một cuốn sách, được kết quả: $2,3\ cm$; $2,4\ cm$; $2,5\ cm$; $2,4\ cm$. Giá trị trung bình của các lần đo là
\choice
{$2,3\ cm$}
{\True $2,4\ cm$}
{$2,5\ cm$}
{$2,2\ cm$}
\haidong{2}
\loigiai{
}
\haidong{7}\end{ex}
\begin{ex}%book
Một người đo chiều dài một cuốn sách $\ell=22 \pm 1\ cm$. Sai số tỉ đối của phép đo chiều dài quyển sách là
\choice
{$4\%$}
{$5\%$}
{\True $4,5\%$}
{$1,5\%$}
\haidong{2}
\loigiai{
}
\haidong{10}\end{ex}
\begin{ex}%book %[1P1BB-2][Loai1]
Khi đo gia tốc rơi tự do, một học sinh tính được $\overline{g}=9,786345\ m/s^2$, $\Delta g=0,025479\ m/s^2$ thì kết quả được ghi như thế nào?
\choice
{$9,78 \pm 0,0254$}
{$9,78 \pm 0,025$}
{\True $9,786 \pm 0,025$}
{$9,786 \pm 0,03$}
\haidong{5}
\loigiai{
}
\haidong{10}\end{ex}
\begin{ex}%book
Đo bề dày $\ell$ của một mặt bàn $4$ lần, ta được kết quả: $2,3\ cm$; $2,4\ cm$; $2,5\ cm$; $2,4\ cm$. Lấy sai số dụng cụ là $0,05\ cm$. Kết quả của phép đo là
\choice
{$\ell=2,40 \pm 0,05\ cm$}
{$\ell=2,50 \pm 0,05\ cm$}
{$\ell=2,3 \pm 0,1\ cm$}
{\True $\ell=2,4 \pm 0,1\ cm$}
\haidong{5}
\loigiai{
}
\haidong{10}\end{ex}
\begin{ex}%book
Có hai điện trở: $(3,0 \pm 0,1)\ \Omega$ và $(6,0 \pm 0,3)\ \Omega$. Phép đo điện trở tương đương của hai điện trở mắc nối tiếp sẽ có sai số tỉ đối bằng
\choice
{$1,1\%$}
{$2,2\%$}
{$3,3\%$}
{\True $4,4\%$}
\haidong{5}
\loigiai{

\begin{itemize}
\item $\overline{R_{n t}}=\overline{R_1}+\overline{R_2}=3,0+6,0=9,0\ \Omega$

\item $\Delta R_{n t}=\Delta R_1+\Delta R_1=0,1+0.3=0,4\ \Omega$

\item Sai số tỉ đối: $\delta\left(R_{n t}\right)=\dfrac{\Delta R_{n t}}{R_{m t}} .  100\%=\dfrac{0,4}{9,0} .  100\%=4,4\%$
\end{itemize}


}
\haidong{10}\end{ex}
\subsubsection*{PHẦN 2. Thí sinh trả lời từ câu 1 đến câu 4. Trong mỗi ý a), b), c), d) ở mỗi câu, thí sinh chọn đúng hoặc sai.}
\setcounter{ex}{0}
\begin{ex}
Về sai số trong phép đo vật lí.
\choiceTF[t]
{\True Sai số tuyệt đối của lần đo thứ $i$ được xác định bởi công thức $\Delta A_i = \left| \overline{A} - A_i \right|$}
{\True Thực hiện phép đo nhiều lần rồi lấy trung bình là cách giảm sai số ngẫu nhiên}
{Sai số hệ thống có thể loại bỏ hoàn toàn bằng cách thực hiện nhiều phép đo}
{Sai số tỉ đối có đơn vị là phần trăm hoặc không có đơn vị tùy theo trường hợp}
\loigiai{
\begin{itemchoice}
\itemch
\itemch
\itemch
\itemch
\end{itemchoice}
}
\end{ex}
\begin{ex}
Liên hệ giữa sai số và độ chính xác của phép đo.
\choiceTF[t]
{\True Sai số tỉ đối càng nhỏ thì phép đo càng chính xác}
{\True Để giảm sai số hệ thống, có thể hiệu chỉnh và sử dụng thiết bị chính xác hơn}
{Sai số ngẫu nhiên có thể loại bỏ hoàn toàn bằng cách thay đổi thiết bị đo}
{Khi sai số tuyệt đối bằng 0 thì phép đo được coi là hoàn hảo và không có sai số nào khác}
\loigiai{
\begin{itemchoice}
\itemch
\itemch
\itemch
\itemch
\end{itemchoice}
}
\end{ex}
\begin{ex}
Về cách xác định số chữ số có nghĩa và sai số trong phép đo.
\choiceTF[t]
{Số 6,07 có 3 chữ số có nghĩa vì các chữ số 6, 0, 7 đều được ghi rõ}
{\True Nếu sai số tuyệt đối là $0,03\ cm$ thì kết quả đo phải được viết chính xác đến hàng phần trăm}
{Số chữ số có nghĩa luôn bằng số chữ số sau dấu phẩy}
{Một kết quả đo có sai số là $0,04\ cm$ thì số đo $12,0\ cm$ là không phù hợp}
\loigiai{
\begin{itemchoice}
\itemch
\itemch
\itemch
\itemch
\end{itemchoice}
}
\end{ex}
\begin{ex}
Về cách xác định và giảm sai số trong thực hành đo.
\choiceTF[t]
{\True Trong phép đo $v = \dfrac{s}{t}$, sai số tỉ đối được xác định bởi $\delta v = \delta s + \delta t$}
{\True Sai số dụng cụ được xác định bằng một nửa độ chia nhỏ nhất của thước hoặc đồng hồ}
{Sai số của phép đo chiều dài không phụ thuộc vào cách đặt mắt khi đọc kết quả}
{Đo khối lượng riêng là phép đo trực tiếp vì dùng cân và bình chia độ}
\loigiai{
\begin{itemchoice}
\itemch
\itemch
\itemch
\itemch
\end{itemchoice}
}
\end{ex}
\subsubsection*{PHẦN 3. Thí sinh trả lời từ câu 1 đến câu 6.}
\setcounter{ex}{0}
\begin{ex}%Tailieuchuan
Một học sinh đo chiều dài 5 lần và thu được các giá trị: $19,8\ cm$, $20,0\ cm$, $19,9\ cm$, $20,1\ cm$, $19,8\ cm$. Giá trị trung bình của chiều dài là bao nhiêu $cm$ (làm tròn kết quả đến chữ số hàng phần mười)?
\shortans[oly]{19,9}
\loigiai{
$\overline{l} = \dfrac{19,8 + 20,0 + 19,9 + 20,1 + 19,8}{5} = 99,6 : 5 = 19,92 \approx 19,9\ cm$
}
\end{ex}

\begin{ex}%Tailieuchuan
Giá trị trung bình của các phép đo: $10,2\ s$, $10,0\ s$, $10,3\ s$, $10,1\ s$ là bao nhiêu giây (làm tròn kết quả đến chữ số hàng phần mười)?
\shortans[oly]{10,2}
\loigiai{
$\overline{x} = \dfrac{10,2 + 10,0 + 10,3 + 10,1}{4} = \dfrac{40,6}{4} = 10,15 \approx 10,2$
}
\end{ex}

\begin{ex}%Tailieuchuan
Đo chiều dài bằng thước có độ chia nhỏ nhất $1\ mm$. Sai số tuyệt đối của phép đo là bao nhiêu $mm$ nếu lấy bằng một nửa độ chia nhỏ nhất?
\shortans[oly]{0,5}
\loigiai{
Sai số tuyệt đối thường lấy bằng $\dfrac{1}{2}$ độ chia nhỏ nhất $\Rightarrow \Delta l = 0,5\ mm$
}
\end{ex}

\begin{ex}%Tailieuchuan
Đo điện áp bằng vôn kế có độ chia nhỏ nhất $0,2\ V$. Sai số tuyệt đối là bao nhiêu vôn nếu lấy bằng một nửa độ chia nhỏ nhất?
\shortans[oly]{0,1}
\loigiai{
$\Delta U = \dfrac{1}{2} .  0,2 = 0,1\ V$
}
\end{ex}

\begin{ex}%Tailieuchuan
Một học sinh đo chiều dài $\ell = 25,4 \pm 0,1\ cm$. Sai số tỉ đối là bao nhiêu $\%$ (làm tròn kết quả đến chữ số hàng phần mười)?
\shortans[oly]{0,4}
\loigiai{
$\delta = \dfrac{0,1}{25,4} .  100\% \approx 0,39\% \approx 0,4\%$
}
\end{ex}

\begin{ex}%Tailieuchuan
Một học sinh đo thời gian rơi tự do được $t = 0,48 \pm 0,02\ s$. Sai số tỉ đối bao nhiêu $\%$  (làm tròn kết quả đến chữ số hàng phần mười)?
\shortans[oly]{4,2}
\loigiai{
$\delta = \dfrac{0,02}{0,48} .  100\% \approx 4,17\% \approx 4,2\%$
}
\end{ex}









\TieudeT{PHONG TỎA VẬT LÍ}
\subsubsection*{PHẦN I. Thí sinh trả lời câu hỏi từ câu 1 đến câu 18. Mỗi câu hỏi thí sinh chỉ chọn một phương án}
\setcounter{ex}{0}
\begin{ex}%book
Hai đơn vị nào sau đây gồm có một đơn vị cơ bản và một đơn vị dẫn xuất trong hệ SI?
\choice
{Kelvin, mol}
{Mét, kilogam}
{\True Mol, jun}
{Giây, ampe}	
\loigiai{
}
\end{ex}
\begin{ex}%book
Đơn vị đo khối lượng trong hệ thống đo lường SI là
\choice
{tấn}
{gam}
{\True kilôgam}
{miligam}
\loigiai{

}
\end{ex}
\begin{ex}%book
Phép đo đại lượng nào sau đây là phép đo trực tiếp?
\choice
{\True Đo thời gian chuyển động của một vật}
{Tốc độ trung bình}
{Đo diện tích}
{Đo gia tốc}	
\loigiai{

}
\end{ex}
\begin{ex}%book
Phép đo đại lượng nào sau đây là phép đo gián tiếp?
\choice
{Quãng đường}
{\True Diện tích}
{Khối lượng}
{Thời gian}
\loigiai{
}
\end{ex}
\begin{ex}%book
Để xác định thành tích chạy của vận động viên điền kinh người ta sử dụng loại đồng hồ nào sau đây?
\choice
{Đồng hồ đeo tay}
{\True Đồng hồ bấm giây}
{Đồng hồ quả lắc}
{Đồng hồ hẹn giờ}	
\loigiai{

}
\end{ex}
\begin{ex}%book %[1P1YB-2][Loai1]
Theo quy ước, số $42,50$ có bao nhiêu chữ số có nghĩa?
\choice
{$1$}
{$2$}
{\True $4$}
{$3$}
\loigiai{
Theo quy ước, số 42,50 có 4 chữ số có nghĩa.
}
\end{ex}
\begin{ex}%book
Xét theo nguyên nhân thì sai số của phép đo được chia thành mấy loại?
\choice
{$1$}
{$3$}
{\True $2$}
{$4$}	
\loigiai{

}
\end{ex}
\begin{ex}%book
Sai số dụng cụ thường lấy bằng
\choice
{\True nửa hoặc một độ chia nhỏ nhất trên dụng cụ đo}
{nửa hoặc một phần tư độ chia nhỏ nhất trên dụng cụ đo}
{nửa hoặc hai độ chia nhỏ nhất trên dụng cụ đo}
{một hoặc hai độ chia nhỏ nhất trên dụng cụ đo}	
\loigiai{
}
\end{ex}
\begin{ex}%book
Sai số tuyệt đối trung bình của n lần đo được tính theo công thức:
\choice
{$A=\overline{A} \pm \Delta A$}
{$\delta A=\dfrac{\Delta A}{\overline{A}} .  100\%$}
{$\overline{A}=\dfrac{A_1+A_2+\ldots+A_{n}}{n}$}
{\True $\overline{\Delta A}=\dfrac{\Delta A_1+\Delta A_2+\ldots+\Delta A_{n}}{n}$}

\loigiai{
}
\end{ex}
\begin{ex}%book
Sai số phép đo gồm
\choice
{sai số hệ thống và sai số dụng cụ}
{sai số chủ quan va sai số khách quan}
{\True sai số hệ thống và sai số ngẫu nhiên}
{sai số điều kiện thí nghiệm và sai số ngẫu nhiên}

\loigiai{
}
\end{ex}
\begin{ex}%book
Phát biểu nào sau đây \textbf{không} đúng về sai số tỉ đối (còn được gọi là sai số tương đối)?
\choice
{Sai số tỉ đối (còn được gọi là sai số tương đối) càng nhỏ thì phép đo càng chính xác}
{\True Sai số tỉ đối (còn được gọi là sai số tương đối) càng lớn thì phép đo càng chính xác}
{Sai số tỉ đối (còn được gọi là sai số tương đối) là tỉ số giữa sai số tuyệt đối và giá trị trung bình}
{Sai số tỉ đối (còn được gọi là sai số tương đối) được tính bằng công thức $\delta x=\dfrac{\Delta x}{\overline{x}} .  100\%$}

\loigiai{
}
\end{ex}
\begin{ex}%book
Kết quả đo đại lượng A được viết dưới dạng $A=\overline{A} \pm \Delta A$. Giá trị thực của đại lượng cần đo A nằm trong khoảng
\choice
{từ $-\Delta A$ đến $+\Delta A$}
{\True từ $\overline{A}-\Delta A$ đến $\overline{A}+\Delta A$}
{từ $-\overline{A}-\Delta A$ đến $\overline{A}+\Delta A$}
{từ $\overline{A}-2 \Delta A$ đến $\overline{A}+2 \Delta A$}

\loigiai{
}
\end{ex}
\begin{ex}%book
Cho các bước khi thực hiện đo nhiệt độ của một vật gồm:
\begin{enumEX}[\faEdit]{2}
\item[(1)] Thực hiện phép đo nhiệt độ.
\item[(2)] Đọc và ghi kết quả đo.
\item[(3)] Lựa chọn nhiệt kế phù hợp.
\item[(4)] Uớc lượng nhiệt độ của vật.
\item[(5)] Hiệu chỉnh nhiệt kế.
\end{enumEX}
Thứ tự đúng khi thực hiện phép đo nhiệt độ là
\choice
{$(1)$, $(2)$, $(3)$, $(4)$, $(5)$}
{\True $(4)$, $(3)$, $(5)$, $(1)$, $(2)$}
{$(2)$, $(4)$, $(3)$, $(1)$, $(5)$}
{$(3)$, $(4)$, $(1)$, $(2)$, $(5)$}

\loigiai{
%		
}
\end{ex}
\begin{ex}%book
Dùng một đồng hồ đo thời gian có độ chia nhỏ nhất $0,001\ s$ để đo thời gian rơi tự do của một vật. Xác định sai số dụng cụ của phép đo \textit{có thể} là
\choice
{\True $0,0005\ s$}
{$0,002\ s$}
{$0,003\ s$}
{$0,004\ s$}
\haidong{3}	
\loigiai{
}
\haidong{7}\end{ex}
\begin{ex}%book
Kết quả đo gia tốc rơi tự do được viết dưới dạng: $g=9,78 \pm 0,44\ m/s^2$. Sai số tỉ đối của phép đo xấp xỉ là
\choice
{$4,0\%$}
{$5,0\%$}
{\True $4,5\%$}
{$3,5\%$}
\haidong{2}
\loigiai{
}
\haidong{10}
\end{ex}
\begin{ex}%book
Thể tích của hai vật đo được bằng $V_1=(10,20 \pm 0,02)\ cm^3$ và $V_2=(6,40 \pm 0,01)\ cm^3$. Tổng thể tích của hai vật trên sẽ có giá trị bằng
\choice
{$(17,00 \pm 0,01)\ cm^3$}
{$(16,60 \pm 0,02)\ cm^3$}
{$(16,60 \pm 0,01)\ cm^3$}
{\True $(16,60 \pm 0,03)\ cm^3$}
\haidong{5}
\loigiai{

$
\begin{aligned}
& V=V_1+V_2 \Rightarrow \overline{V}=\overline{V}_1+\overline{V}_2=10,2+6,4=16,6\ cm^3 \\
& \Delta V=\Delta V_1+\Delta V_2=0,02+0,01=0,03\ cm^3 \\
& \Rightarrow V=(16,60 \pm 0,03)\ cm^3
\end{aligned}
$

}
\haidong{10}\end{ex}
\begin{ex}%book %[1P1BB-2][Loai1] 
\immini[thm]{Dùng thước kẹp có độ chia nhỏ nhất $0,1\ mm$ để đo $5$ lần chiều cao $h$ của một trụ thép. Cho kết quả như trong bảng dưới. Lấy sai số dụng cụ là một độ chia nhỏ nhất. Kết quả phép đo $h$ của trụ thép là
\choice
{$h=(19,9\pm 0,18)\ mm$}
{\True $h=(19,9\pm 0,2)\ mm$}
{$h=(19,8\pm 0,2)\ mm$}
{$h=(19,9\pm 0,1)\ mm$}}{\begin{tabular}{|c|c|}
\hline 
Lần đo&h (mm)  \\ 
\hline 
1&19,9  \\ 
\hline 
2&19,8  \\ 
\hline 
3&20,0  \\ 
\hline 
4&19,7  \\ 
\hline 
5& 19,9 \\ 
\hline 
\end{tabular}}
\haidong{4}
\loigiai{
\begin{itemize}
\item Giá trị trung bình: $\overline{h}=\dfrac{19,9+19,8+20+19,9+19,7}{5}=19,86\ mm$. 
\item Sai số tuyệt đối của mỗi phép đo lần lượt là: $0,04;\ 0,06;\ 0,14;\ 0,16;\ 0,04$.
\item Sai số ngẫu nhiên: $\overline{\Delta h}=0,088\ mm$.
\item Sai số dụng cụ: $\Delta  h_{DC}=0,1\ mm\Rightarrow\Delta h=\Delta  h_{DC}+\overline{\Delta h}=0,188\ mm\approx 0,2\ mm$.
\item Kết quả phép đo: $h=(19,9\pm 0,2)\ mm$.
\end{itemize}
}
\haidong{10}\end{ex}

\begin{ex}%book
Biết rằng đại lượng vật lí $(A)$ được đo gián tiếp thông qua hai đại lượng vật lí $(B)$ và $(C)$ bằng công thức toán học $A=\dfrac{B}{C^2}$. Kết quả đo đại lượng $(B)$ và $(C)$ lần lượt là $B=55,35 \pm 0,44$; $C=2,86 \pm 0,12$. Kết quả của đại lượng $(A)$ được ghi là
\choice
{$A=18,30 \pm 0,62$}
{$A=6,77 \pm 0,34$}
{$A=18,30 \pm 0,34$}
{\True $A=6,77 \pm 0,62$}
\haidong{9}
\loigiai{
\begin{itemize}
\item $\overline{A}=6,77$
\item $\delta A=\delta B+2\delta C\Rightarrow \dfrac{\Delta A}{\overline{A}}=\dfrac{\Delta B}{\overline{B}}+2.\dfrac{\Delta C}{\overline{C}}\Rightarrow \dfrac{\Delta A}{6,77}=\dfrac{0,44}{55,35}+2.\dfrac{0,12}{2,86}\Rightarrow \Delta A=0,62$
\end{itemize}
}
\haidong{10}\end{ex}










\subsubsection*{PHẦN 2. Thí sinh trả lời từ câu 1 đến câu 4. Trong mỗi ý a), b), c), d) ở mỗi câu, thí sinh chọn đúng hoặc sai.}
\setcounter{ex}{0}
\begin{ex}%Sai số phép đo - Đúng Sai
Trong quá trình đo đạc vật lí, người ta thường lặp lại nhiều lần để giảm sai số.
\choiceTF[t]
{\True Giá trị trung bình của nhiều phép đo lặp lại giúp giảm ảnh hưởng của sai số ngẫu nhiên}
{Càng đo nhiều lần thì sai số hệ thống sẽ tự triệt tiêu}
{\True Dụng cụ đo có độ chia càng nhỏ thì sai số dụng cụ càng nhỏ}
{Kết quả đo phải luôn trùng khớp với giá trị thực của đại lượng}
\loigiai{
\begin{itemchoice}
\itemch Đúng. Trung bình cộng nhiều phép đo làm giảm ảnh hưởng của sai số ngẫu nhiên.
\itemch Sai. Sai số hệ thống mang tính lặp lại có quy luật nên không bị triệt tiêu bởi việc lặp lại phép đo.
\itemch Đúng. Độ chia nhỏ hơn giúp tăng độ chính xác, giảm sai số dụng cụ.
\itemch Sai. Kết quả đo luôn chứa sai số, khó trùng khớp tuyệt đối với giá trị thực.
\end{itemchoice}
}
\end{ex}

\begin{ex}%Sai số phép đo - Đúng Sai
Phân biệt giữa sai số tuyệt đối và sai số tỉ đối trong đo lường vật lí.
\choiceTF[t]
{\True Sai số tuyệt đối có cùng đơn vị với đại lượng cần đo}
{Sai số tỉ đối được tính bằng hiệu giữa các lần đo chia cho giá trị nhỏ nhất}
{\True Sai số tỉ đối không có đơn vị}
{Sai số tuyệt đối càng nhỏ thì sai số tỉ đối luôn nhỏ}
\loigiai{
\begin{itemchoice}
\itemch Đúng. Sai số tuyệt đối có cùng đơn vị với đại lượng đo.
\itemch Sai. Sai số tỉ đối là tỉ lệ giữa sai số tuyệt đối và giá trị đo trung bình.
\itemch Đúng. Vì là tỉ số nên sai số tỉ đối không mang đơn vị.
\itemch Sai. Nếu giá trị đo rất nhỏ thì sai số tỉ đối vẫn có thể lớn, dù sai số tuyệt đối nhỏ.
\end{itemchoice}
}
\end{ex}

\begin{ex}%Sai số phép đo - Đúng Sai
Ảnh hưởng của yếu tố con người và điều kiện môi trường đến sai số đo.
\choiceTF[t]
{Mắt người đọc không cùng hướng vuông góc với thước sẽ gây ra sai số hệ thống}
{\True Gió thổi mạnh làm sai lệch khi đo thời gian rơi tự do là ví dụ về sai số ngẫu nhiên}
{Người đọc sai đơn vị trên thước là sai số ngẫu nhiên}
{\True Sai số do độ trễ của đồng hồ bấm giây là sai số hệ thống}
\loigiai{
\begin{itemchoice}
\itemch Sai. Sai lệch do góc nhìn sai gây ra sai số ngẫu nhiên, không mang tính quy luật.
\itemch Đúng. Điều kiện môi trường thay đổi bất thường gây sai số ngẫu nhiên.
\itemch Sai. Đọc sai đơn vị là sai sót cá nhân, không được xếp vào sai số ngẫu nhiên.
\itemch Đúng. Độ trễ của dụng cụ là sai số có tính lặp lại, nên là sai số hệ thống.
\end{itemchoice}
}
\end{ex}

\begin{ex}%Sai số phép đo - Đúng Sai
Quy tắc ghi kết quả đo và số chữ số có nghĩa.
\choiceTF[t]
{Kết quả đo chỉ nên ghi dưới dạng $A = \overline{A} \pm \Delta A$}
{Nếu sai số tuyệt đối là $0,05$ thì kết quả có thể ghi là $12,347 \pm 0,05$}
{\True Khi nhân chia hai đại lượng đo, sai số tỉ đối phải cộng vào nhau}
{\True Số chữ số có nghĩa của kết quả đo phải phù hợp với độ chính xác của phép đo}
\loigiai{
\begin{itemchoice}
\itemch Đúng. Đây là cách ghi kết quả phép đo chuẩn.
\itemch Sai. Phần thập phân của kết quả và sai số phải cùng bậc (số sau dấu phẩy phải tương ứng).
\itemch Đúng. Khi nhân hoặc chia các đại lượng, ta cộng sai số tỉ đối.
\itemch Đúng. Không được ghi quá nhiều chữ số có nghĩa vượt quá độ chính xác của phép đo.
\end{itemchoice}
}
\end{ex}


\subsubsection*{PHẦN 3. Thí sinh trả lời từ câu 1 đến câu 6.}
\setcounter{ex}{0}
\begin{ex}%Tailieuchuan
Thực hiện $3$ lần đo khối lượng: $120,0\ g$, $119,8\ g$, $120,2\ g$. Giá trị trung bình là bao nhiêu gam (làm tròn kết quả đến chữ số hàng phần mười)?
\shortans[oly]{120}
\loigiai{
$\overline{m} = \dfrac{120,0 + 119,8 + 120,2}{3} = \dfrac{360,0}{3} = 120,0\ g$
}
\end{ex}

\begin{ex}%Tailieuchuan
Kết quả đo: $x = 8,43 \pm 0,02\ m$. Giá trị lớn nhất của phép đo có thể là bao nhiêu mét?
\shortans[oly]{8,45}
\loigiai{
$x_{\text{max}} = 8,43 + 0,02 = 8,45\ m$
}
\end{ex}

\begin{ex}%Tailieuchuan
Kết quả đo: $x = 5,70 \pm 0,03\ m$. Giá trị nhỏ nhất có thể là bao nhiêu mét?
\shortans[oly]{5,67}
\loigiai{
$x_{\text{min}} = 5,70 - 0,03 = 5,67\ m$
}
\end{ex}

\begin{ex}%Tailieuchuan
Một vật có khối lượng $m = 50,0\ g$ và sai số tuyệt đối $0,2\ g$. Sai số tỉ đối là bao nhiêu $\%$?
\shortans[oly]{0,4}
\loigiai{
$\delta = \dfrac{0{,}2}{50{,}0} \cdot 100\% = 0{,}4\%$
}
\end{ex}
\begin{ex}%Tailieuchuan
Dùng thước có độ chia nhỏ nhất $2\ mm$ để đo chiều dài. Sai số dụng cụ là bao nhiêu $mm$ nếu nó được lấy là một nửa độ chia nhỏ nhất?
\shortans[oly]{1}
\loigiai{
Sai số dụng cụ là $\dfrac{2}{2} = 1\ mm = 1\ mm$
}
\end{ex}


\begin{ex}%Tailieuchuan
Một học sinh đo chiều cao $h = 170,0\ cm \pm 0,5\ cm$. Sai số tỉ đối là bao nhiêu $\%$ (làm tròn kết quả đến chữ số hàng phần mười)?
\shortans[oly]{0,3}
\loigiai{
$\delta = \dfrac{0,5}{170,0} .  100\% \approx 0,29\% \approx 0,3\%$
}
\end{ex}